%Author: AA partly adapted from main.tex on the Overleaf
\section{Materials and Methods}
The values of secondary cases in panel (a) of \fref{ls}  were generated by computing the  geometric probability mass function with mean \Ro~at points $0,1,\ldots,30$.
As for the expected infectiousness, we evaluated
the exponential probability density function with mean \Ro~at $300$ equally distanced point in the
 interval $[0,30]$.

To generate panel (b) of \fref{ls}, we followed the Lloyd-Smith's approach and computed the inequality in the activity distribution \cite[SI, Sec. 2.2.5]{lloydsmith2005superspreading}. 
More specifically, we used the relation $F_{\text{trans}}(x)=\frac{1}{\Ro^2}\int_{0}^x ve^{-\frac{v}{\Ro}}dv$ which is the fraction of cases
due to those caused up to $x$ secondary cases.
The fraction of cases due to those caused more than $x$ cases would be $1-F_{\text{trans}}(x)$, which is equal to $(1+\frac{x}{\Ro})e^{-\frac{x}{\Ro}}$.
The population fraction of the individuals infected more than $x$ cases is $e^{-\frac{x}{\Ro}}$.
We used a similar approach to compute the  inequality in the secondary case distribution: We used the relation $F_{\mathrm{trans}}(x)=\frac{1}{\Ro}\sum_{v=0}^x vG(v,\Ro)$ which is the fraction of cases
due to those caused up to $x$ secondary cases.
Here $G(v,\Ro)$ is a geometric distribution with mean \Ro.
The fraction of cases due to those caused more than $x$ cases equals $1-F_{\text{trans}}(x)$. 
The population fraction of these individuals would be $1 - P(x)$, where $P(x)$ is the cumulative distribution function of a geometric distribution with mean \Ro~evaluated at $x$.

To generate panel (a) of \fref{RcAverage}, we used the individual based simulation and computed the number of secondary cases in a population of size $10^4$. 
The simulation was initiated with one case.

We used a simple deterministic SIR model to compute the mean and variance of expected infectiousness \Rx{c}.
First, we scaled time by the mean infectious period, so the resulting SIR model then depends on one parameter: \Ro. The SIR differential equations were numerically solved for the proportions of susceptible $x(t)$ and infectious $y(t)$ at each time point $t$. 
The time interval for integration was set to $[0,100]$, well after the outbreak died out.
We then used the \textbf{R} function \fun{approxfun} to construct a function,  ${X}(t)$, that linearly interpolates the time-series $(t,\Ro x(t))$.
We associated a cohort to each of the first $60\%$ time points in the time-series $t$ (We used a $60\%$ cutoff to ensure cohorts had almost recovered by the end of the simulation period) and calculated the cohort-specific mean  and variance.
More specifically, for the cohort infected at time point $\tau$, the
mean, $\mu(\tau)$, and variance $\sigma^2(\tau)$ of \Rx{c} read as:
\begin{subequations} \label{eq_statistics}
	\begin{align}
 		\mu(\tau) &=  \int_{t>\tau} f(t-\tau) \int_{\tau}^{t} \Ro x(s) \ds \dt, \\
 		\E[\Rx{c}^2(\tau)] &= \int_{t>\tau} f(t-\tau) \left(\int_{\tau}^{t} \Ro x(s) \ds \right)^2\dt,\\
  		\sigma^2(\tau) &=  \E[\Rx{c}^2(\tau)] - \mu^2(\tau).
  	\end{align}
\end{subequations}
To calculate these quantities for the cohort infected at $\tau$, the interpolating function $X(t)$ was integrated from $\tau$ to $t$, which is the case reproductive potential associated with the cohort fraction recovered at time point $t$.
The mean \Rx{c} for each cohort $\mu(\tau)$ was then the integral of the case reproductive potential weighted by the infectious period density function.
We used the same approach to compute $\E[\Rx{c}^2(\tau)]$ and  variance of \Rx{c} for each cohort ($\sigma^2(\tau)$).
In \fref{Cohort}, the middle panel (panel b) was generated using this approach.

We obtained the mean of expected infectiousness \Rx{c} by integrating the cohort-specific mean $\mu(\tau)$ against the incidence, $i(\tau) = \Ro x(\tau)y(\tau)$, and normalizing the result by the final size, $\int i(t) \dt$.
To compute the between cohort variance, we integrated the squared cohort mean $\mu^2(\tau)$ weighted by the incidence $i(\tau)$ and divided it  by the final size.
The result minus the squared mean of \Rx{c} equaled the between cohort variance.
The within-cohort variance  was computed by taking integral over the cohorts' variance $\sigma^2(\tau)$ weighted by the incidence.
The result was also normalized by the final size.
The total variance was the sum of the between- and within-variances. 
We also calculated the total variance independently; 
  we integrated $\E[\Rx{c}^2(\tau)]$ weighted by the incidence and divided it by the final size. 
 The result minus the squared mean of \Rx{c}~yielded the total variance.
Both approaches produced the same value for the total variance.


We used deterministic simulations to compute the y-axis values in the panels b and c of \Fref{Naive}.
Specifically, for each time point $t_f$, we considered the cohorts infected up to time $t_f$.
We then computed cohort statistics by calculating the cases generated by these cohorts up to time $t_f$.
For example, for a cohort infected at time point $\tau < t_f$, the mean, $	\mu(\tau,t_f)$, and variance of \Rx{c}, 	$\sigma^2(\tau, t_f)$, read as:
\begin{subequations} \label{eq_naive}
	\begin{align}
	\mu(\tau,t_f) &=  \int_{t=\tau}^{t_f} f(t-\tau)dt \int_{\tau}^{t} \Ro x(s) ds + F(t_f - \tau)\int_\tau^{t_f}\Ro x(s)ds,  \\
	\E[\Rx{c}^2(\tau,t_f)] \!&=\!\int_{t=\tau}^{t_f} \!f(t-\tau)dt \left(\int_{\tau}^{t} \!\Ro x(s) ds \right)^2 \nonumber\\
	 & \quad + F(t_f-\tau)\!\left(\int_{\tau}^{t_f} \Ro x(s) ds \right)^2,  \\
	\sigma^2(\tau, t_f) &=  \E[\Rx{c}^2(\tau,t_f)] - \mu^2(\tau,t_f),
\end{align}
\end{subequations}
where $F(\cdot)$ is the survival function associated with the  infectious period density $f(\cdot)$.
The mean and variance of \Rx{c} are then 
\begin{subequations} \label{eq_statistics_naive}
	\begin{align}
	\mu(t_f) &= \frac{\int_0^{t_f}i(\tau)\mu(\tau,t_f)d\tau}{\int_0^{t_f}i(t)dt},\\
	\sigma^2(t_f) &= \frac{\int_0^{t_f}i(\tau)\E[\Rx{c}^2(\tau,t_f)]d\tau}{\int_0^{t_f}i(t)dt} - \mu^2(t_f).
	\end{align}
\end{subequations}
The between-cohort component of $\kappa$ at time point $t_f$ is
$	\kappa_{\rm bet}(t_f) = \sigma^2_{\rm bet}(t_f)/\mu^2(t_f), $
where 
\begin{align*}
	\sigma^2_{\rm bet}(t_f) = \frac{\int_0^{t_f} i(\tau)\mu(\tau,t_f)^2}{\int_0^{t_f}i(t)dt} - \mu^2(t_f).
\end{align*}
o generate panels (b) and (c) of \Fref{timeCutoff}, as with the corresponding panels of \Fref{Naive}, at each time point $t_f$, we considered only cohorts infected by $t_f$.
Instead, we counted the total cases caused by these cohorts during the outbreak.
To do so, we simply substituted $\mu(\tau)$ and $\E[\Rx{c}^2(\tau)]$ for $\mu(\tau,t_f)$ and $\E[\Rx{c}^2(\tau,t_f)]$ in Equation \eqref{eq_statistics_naive}, where
$\mu(\tau)$ and $\E[\Rx{c}^2(\tau)]$ were provided in Equation \eqref{eq_statistics}.

 All integrations were done using \fun{lsoda} method with the \textbf{R} package \pkg{deSolve}.
  All simulations were carried out using \textbf{R} \texttt{4.5.2}.
    Code for all numerical simulations is housed at: \url{https://github.com/dushoff/kappaCode}.   
We solved all integrals across a range of values for $\Ro$, using the starting values $y_0 = 10^{-9}; x_0 = 1-10^{-9}$  to represent the limiting case in which there are no exogenous cases.
%, and therefore the mean case reproduction number over the course of the epidemic must be 1.
 In building these simulations, we used a range of time step sizes, noting convergence towards known and conjectured values (e.g., epidemic final size, mean case reproduction number, variance in case reproduction number) as resolution increased. 


